\section{Falsifiable predictions}

The loop hypothesis makes a concrete, falsifiable commitment in the form of a double dissociation. The two streams carry functionally distinct information---the priority stream the grounded change-evidence that drives detection, the value stream the cue-value and reliability context that modulates it---so selectively disabling each should produce opposite, separable behavioural signatures. Lesioning the value circuit should degrade value and context modulation while sparing change detection; lesioning the priority node should collapse detection itself. We state the prediction at three levels of measurement, each independently testable, and tie each to a corresponding manipulation in the model and in the animal.

\subsection{Behavioural level}

The coarsest test is the one the present comparison already exemplifies and which can be reproduced as an intervention rather than an architectural ablation. Two behavioural laws follow from the loop hypothesis. First, \emph{coupling is required}: disrupting the shared recurrent state---so that the actor and critic no longer occupy a common representation---should abolish learning of the declare response, reproducing the always-wait collapse seen in the independent variant (hit rate $0.000$ versus the working model's $0.794$). Second, \emph{the grounded stream must command the action}: routing the policy read-out from the value stream rather than the priority stream should likewise collapse detection, as in the swapped variant. Both are categorical, all-or-nothing effects---a hit-rate gap of roughly $79$ points from a single structural edit---rather than graded performance changes, which makes them stringent tests: the hypothesis forbids a working model that violates either condition.

The corresponding animal prediction is a dissociation between two manipulations. Inactivating the priority/commitment node---the superior colliculus---should collapse covert change detection, mirroring the loss of the grounded actor input; this is consistent with the established finding that SC inactivation produces a large covert-attention deficit \citep{zenon2012attention,lovejoy2010inactivation,bollimunta2018fefsc}, and the prediction sharpens it to the orientation-change-detection setting and to optogenetically precise, reversible perturbation \citep{krauzlis2018selective,katz2025optogenetic}. By contrast, a manipulation confined to the value circuit---the dopaminergic teaching system or its ventral-striatal target---should leave the animal's ability to detect a change broadly intact while degrading how strongly cue value and cue reliability modulate behaviour. In the model the working variant exhibits exactly the modulation that such a value lesion should target: a graded, reliability-scaled validity benefit, with per-validity hit rates rising monotonically from $0.680$ at validity $0.25$ through $0.697$ and $0.750$ to $0.753$ at validity $1.00$, and, at matched change magnitude, a valid-over-invalid advantage that scales with reliability (for example at a $14^{\circ}$ change, $0.748$ valid versus $0.656$ invalid at validity $0.25$, widening to $0.776$ versus $0.568$ at validity $1.00$). The double dissociation predicts that a priority/SC lesion flattens detection regardless of these contextual variables, whereas a value lesion flattens the reliability-scaled benefit while detection of large, salient changes survives.

\subsection{Representational level}

The streams should be \emph{decodably} specialised. The priority stream's output $Z_{\mathrm{pri}}$, whose residual preserves the raw image, should carry the change-evidence: a decoder trained on $Z_{\mathrm{pri}}$ should read out the presence and location of the orientation change far better than one trained on the value stream's output $Z_{\mathrm{val}}$, whose residual is the deep memory and which carries little live visual content. Conversely, the value/context variables---cue value (reward magnitude) and cue reliability (validity)---should decode far better from $Z_{\mathrm{val}}$ than from $Z_{\mathrm{pri}}$. This is a within-model prediction that is directly verifiable from the activations and that grounds the lesion logic: it asserts that the two streams are not redundant copies but information-segregated channels, so that disabling one removes a specific, identifiable variable. The animal analogue is the expectation that grounded sensory and commitment areas (sensory cortex, FEF-visual, SC) carry the discriminative change-evidence \citep{thompson2005neuronal,gold2007neural,katz2016dissociated}, whereas the striatal/value system carries the reward and probability context on a partly shared substrate \citep{platt1999neural,sugrue2004matching}.

\subsection{Causal-perturbation level}

The strongest test is causal and graded rather than ablative. Perturbing the priority attention computation---scaling or clamping the grounded stream that feeds the actor---should shift behaviour and the detection threshold: a weakened priority signal should raise the effective change magnitude needed to declare, lengthen reaction times (which in the intact model fall monotonically with change magnitude, from a median of $10$ frames at $2^{\circ}$ to $1$ frame at $44^{\circ}$), and ultimately abolish detection as the perturbation deepens. Perturbing the value attention computation---clamping the stream that feeds the critic---should instead shift the critic's value estimates and the cue-value and reliability modulation while leaving the detection threshold and reaction-time structure for large changes essentially unchanged. This is the model-internal realisation of the double dissociation: priority perturbation moves the \emph{behavioural threshold}, value perturbation moves the \emph{critic's valuation}, and the two effects are separable.

The matched animal experiment is graded optogenetic or microstimulation control of the superior colliculus versus a manipulation of the dopaminergic value circuit. Subthreshold collicular and frontal manipulations are known to shift covert attentional performance in a graded fashion \citep{moore2001control,moore2004microstimulation,stine2023differential,herman2020attention}, and recent optogenetic SC perturbation provides the precise, reversible tool the prediction requires \citep{katz2025optogenetic,krauzlis2018selective}: graded SC suppression should produce a graded rise in detection threshold up to collapse, whereas a value-circuit manipulation of comparable behavioural scale should leave detection of salient changes intact while selectively blunting the reliability-scaled and value-scaled modulation of choice. A finding that an SC perturbation spared detection, or that a value-circuit perturbation collapsed detection of large changes, would falsify the loop hypothesis as stated.
